%!TEX root = ../../thesis-trier.tex

Here, we investigate the constrained synchronization problem for monotonic input automata.
When not said otherwise, by $\le$ we denote the linear order given on the states
of some monotonic automaton.

\Cref{constr_sync:lem:monotonic_min_max_sync} first states 
a characterization of synchronizing monotonic automata.
Then, in~\cref{constr_sync:prop:monotonic_bound} 
we give a quadratic upper bound for a synchronizing word in a given constraint language.
\Cref{constr_sync:thm:monotonic_in_NL}
states that the constrained synchronization problem
for monotonic input automata is in~\NL, 
even when we have no access to the linear order on the states.
%where we make no assumption on the encoding of the linear order, or that it is given with the automaton
%apart from mere existence of this order. 

\begin{lemma}
	\label{constr_sync:lem:monotonic_min_max_sync}
	Let $\semiautomat{A}{Q}{\Sigma}{\delta}$ be a monotonic automaton
	where $s \in Q$ is the minimum in the linear order on the states
	and $t \in Q$ the maximum.
	Then $u \in \Sigma^*$ synchronizes $\mathcal A$
	if and only if $\applytrans{\delta}{s}{u} = \applytrans{\delta}{t}{u}$.
\end{lemma}
\begin{proof}
	If $u \in \Sigma^*$ synchronizes $\mathcal A$,
	then in particular  $\applytrans{\delta}{s}{u} = \applytrans{\delta}{t}{u}$.
	Conversely, for every $q \in Q$ we have $s \le q \le t$
	and so $\applytrans{\delta}{s}{u} \le \applytrans{\delta}{q}{u} \le \applytrans{\delta}{t}{u}$
	for $u \in \Sigma^*$.
	Hence, if $\applytrans{\delta}{s}{u} = \applytrans{\delta}{t}{u}$,
	then we get $\applytrans{\delta}{Q}{u} = \{ \applytrans{\delta}{t}{u} \}$. 
\end{proof}

The following proof uses the pair automaton (defined in Section~\ref{sec:pair_aut}).

\begin{proposition}
	\label{constr_sync:prop:monotonic_bound}
	Let $\automat{B}{P}{\Sigma}{\mu}{p_0}{F}$ be a PDFA. % constraint PDFA.
	If $\semiautomat{A}{Q}{\Sigma}{\delta}$
	is an $n$-state monotonic automaton with a synchronizing word in $L(\mathcal B)$,
	then a shortest word in~$\Syn(\mathcal A) \cap L(\mathcal B)$
	has length at most $\left( \binom{n}{2} + 1 \right) \cdot |P| - 1$.
\end{proposition}
\begin{proof} 
	We can assume $s \ne t$. Consider the pair automaton $\mathcal C = \mathcal A^{\le 2} / {\sim_S}$
	where $S = \{ \{q\} \mid q \in Q \}$.
	Let $s \in Q$ be the minimum and $t \in Q$
	be the maximum state.
	By~\cref{constr_sync:lem:monotonic_min_max_sync} 
	we have $\Syn(\mathcal A) = L(\mathcal C_{\{s,t\}, \{S\}})$.
	By \cref{prelim:eqn:product-intersection} 
	the language $L(\mathcal C_{\{s,t\}, \{S\}}) \cap L(\mathcal B)$
	is accepted by the product PDFA
	of~$\mathcal C$ and~$\mathcal B$
	with start state $(\{s,t\}, p_0)$ and final state set $\{S\}\times F$,
	which has 
	$\left(\binom{n}{2} + 1 \right) \cdot |P|$ states in total.
	So, a shortest path from from $(\{s,t\}, p_0)$
	to a state $(S, p)$ with $p \in F$
	has length at most $\left(\binom{n}{2} + 1 \right) \cdot |P| - 1$.	
\end{proof} 

\begin{theorem}
	\label{constr_sync:thm:monotonic_in_NL}	
	Let $L \subseteq \Sigma^*$ be fixed regular language. The problem to decide whether a given monotonic input automaton $\semiautomat{A}{Q}{\Sigma}{\delta}$ 
	has a synchronizing word in $L$ is in~\NL.
\end{theorem}
\begin{proof}
	Set $n = |Q|$. We can assume $n \ge 2$.
	Let $\automat{B}{P}{\Sigma}{\mu}{p_0}{F}$ be a PDFA
	with $L = L(\mathcal B)$. 
	Loop through distinct pairs of states
	$p, q \in Q$ and nondeterministically guess a word $u \in \Sigma^*$ of length at most $\left( \binom{n}{2} + 1 \right) \cdot |P| - 1$
	letter-by-letter %while updating $\applytrans{\delta}{\{p,q\}}{u}$.
	with $\applytrans{\mu}{p_0}{u} \in F$ and $|\applytrans{\delta}{\{p,q\}}{u}| = 1$.
	%After each iteration for two states $p,q \in Q$, 
	%check whether $\applytrans{\mu}{p_0}{u} \in F$
	%and $|\applytrans{\delta}{\{p,q\}}{u}| = 1$.
	This procedure can run in log-space by using counters (or storage space) with a binary encoding
	for the states and the length of the word.
	Then $\mathcal A$ has a synchronizing word  in $L(\mathcal B)$
	if and only if a word $u \in L(\mathcal B)$ as above
	can be found for each pair of states $p, q \in Q$.
	For if there exists a synchronizing word in~$L(\mathcal B)$,
	then by~\cref{constr_sync:prop:monotonic_bound} it has length at most $\left( \binom{n}{2} + 1 \right) \cdot |P| - 1$.
	Hence, in each iteration for two states, by nondeterminism,
	a word in $L(\mathcal B)$ is guessed that maps both states to a single 
	state.
	Conversely, if each iteration for two states finds a word in $L(\mathcal B)$
	mapping both to a single state,
	then in particular for
	the minimum state $s \in Q$ and the maximum state $t \in Q$
	we find a word $u \in L(\mathcal B)$ such that $\applytrans{\delta}{s}{u} = \applytrans{\delta}{t}{u}$.
	So, by~\cref{constr_sync:lem:monotonic_min_max_sync},
	the word $u \in L(\mathcal B)$ synchronizes~$\mathcal A$. 	
\end{proof} 

However, we do not know whether the above problem for monotonic automata 
is \NL-hard.